\section{Communication and System Architecture}
\label{sec:communication}


This document describes the internal communication flow of \texttt{xeus-haskell}, from the Jupyter frontend down to the MicroHs runtime.
\subsection{System Stack}


The communication stack is organized as follows:
Jupyter Frontend \(\leftrightarrow\) Xeus C++ Library \(\leftrightarrow\) Xeus-Haskell Kernel \(\leftrightarrow\) \texttt{mhs\_repl.cpp} Bridge \(\leftrightarrow\) Repl.lhs Compiler API \(\leftrightarrow\) MicroHs Runtime.


\subsection{Layer Details}


\subsubsection{1. Jupyter and Xeus Protocol Layer}


This layer handles the ZeroMQ-based Jupyter protocol. It manages the heartbeat, shell, and control sockets.
\subsubsection{2. Xeus to Kernel Interpreter Layer}


The \texttt{xinterpreter} implementation hooks into the Xeus kernel lifecycle. It receives execution requests and delegates them to the Haskell backend.
\subsubsection{3. C++ Bridge to Haskell Runtime Layer}


This is the bridge layer between C++ and the Haskell-compiled Compiler API wrapper (\texttt{Repl.lhs}).
The \texttt{mhs\_repl.cpp} bridge performs three primary tasks:
\begin{enumerate}
\item \textbf{Initialize MicroHs Runtime}:
   The \texttt{MicroHsRepl} constructor locates the runtime files (standard libraries) using \texttt{microhs\_runtime\_dir}. After initialization, it evaluates a trivial expression (\texttt{"0"}) to warm up the compiler and cache library compilation.\end{enumerate}


\begin{enumerate}
\item \textbf{Capture Standard Output}:
   To receive output from the Haskell world, we use POSIX pipes to redirect \texttt{stdout}.   - C++ sets up a redirection from \texttt{stdout} to a temporary buffer.   - After the Haskell evaluation completes, the buffer is read.   - For implementation details, see \texttt{capture\_stdout} in the source code.\end{enumerate}


\begin{enumerate}
\item \textbf{Parse Captured Content}:
   The output is parsed to distinguish between plain text and rich content (HTML, LaTeX, etc.) using the internal protocol described below.\end{enumerate}


\subsubsection{4. REPL Engine and MicroHs Compiler Layer}


This layer uses the MicroHs Compiler API to evaluate Haskell code strings in the runtime environment.
\subsection{Rich Content Protocol}


\texttt{xeus-haskell} supports rendering rich media (HTML, images, LaTeX) by embedding control sequences in the standard output stream.
\subsubsection{Protocol Structure}


Rich content is strictly delimited by \textbf{ASCII control characters}:
\begin{minted}{text}
<STX><MIME_TYPE><US><CONTENT><ETX>
\end{minted}


\begin{center}
\begin{tabular}{|l|c|p{8cm}|}
\hline
\textbf{Component} & \textbf{ASCII (Hex)} & \textbf{Description} \\
\hline
\texttt{STX} & \texttt{02} & \textbf{S}tart of \textbf{T}e\textbf{x}t \\
\hline
\texttt{MIME\_TYPE} & - & The target MIME type (e.g., \texttt{text/html}) \\
\hline
\texttt{US} & \texttt{1F} & \textbf{U}nit \textbf{S}eparator \\
\hline
\texttt{CONTENT} & - & The raw data/string to be rendered \\
\hline
\texttt{ETX} & \texttt{03} & \textbf{E}nd of \textbf{T}e\textbf{x}t \\
\hline
\end{tabular}
\end{center}
\subsubsection{Supported Content Types}


\begin{itemize}
\item \texttt{text/plain}: Standard console output.
\item \texttt{text/html}: Rendered as HTML in Jupyter.
\item \texttt{text/latex}: Mathematical equations rendered via MathJax.
\item \texttt{text/markdown}: Formatted text and tables.
\end{itemize}


All content is expected to be \textbf{UTF-8} encoded.
\subsection{Haskell Display API}


The protocol is abstracted in Haskell through the \texttt{Display} typeclass.
\subsubsection{1. Display Typeclass Interface}


To enable rich output for a custom type, implement the \texttt{Display} typeclass:
\begin{minted}{haskell}
import XHaskell.Display

-- | Data structure for the protocol
data DisplayData = DisplayData {
    mimeType :: String,
    content  :: String
}
\end{minted}


\subsubsection{2. Implementation Examples}


\begin{minted}{haskell}
newtype HTMLString = HTMLString String
instance Display HTMLString where
    display (HTMLString s) = DisplayData "text/html" s

newtype LaTeXString = LaTeXString String
instance Display LaTeXString where
    display (LaTeXString s) = DisplayData "text/latex" s

-- Default instance for plain strings
instance Display String where
    display s = DisplayData "text/plain" s
\end{minted}


\subsubsection{3. Protocol Generation}


The \texttt{DisplayData} record implements a \texttt{Show} instance that automatically formats the string into the \texttt{<STX>...<ETX>} protocol.
\begin{minted}{haskell}
putStr $ show (display (HTMLString "<p style='color: red'>Hello Haskell!</p>"))
\end{minted}


\texttt{02} \texttt{text/html} \texttt{1F} \texttt{<p style='color: red'>Hello Haskell!</p>} \texttt{03}
